\section{Mechanistic analysis}
\label{sec:mechanism}

\subsection{How does FR help? (Q2)}
The mechanism is a feedback loop between coverage and the mean-force estimate.
\begin{enumerate}
  \item ABF needs samples \emph{across} the reaction coordinate: by
        \eqref{eq:meanforce}, $\Fpref(x)$ is a conditional average that can only
        be estimated where particles actually visit $x$.
  \item Early in a run the bias is poor and the ensemble is unbalanced ---
        \cref{fig:mechanism}(b) shows the initial occupancy split
        ($\sim\!0.5$ in one basin) relaxing only gradually toward the balanced
        $1/3$--$1/3$--$1/3$ split across left basin, barrier and right basin.
  \item FR compares the empirical marginal $\phat_t$ with the target $q_t$:
        particles in over-represented regions of $x$ receive positive
        score~\eqref{eq:score} and are more likely to be killed, while those in
        under-represented regions receive negative score and are more likely to
        be cloned~\eqref{eq:death}.
  \item This raises the useful coverage / sampling balance of $x$ earlier than
        diffusion alone would. \Cref{fig:mechanism}(a) shows the estimated-FR
        free-energy error sitting below ABF only in the early window, and the
        FR score dispersion $\mathrm{std}_i\,S_i$ in panel (d) decaying as
        $\phat_t$ is pulled toward $q_t$.
  \item Better coverage feeds back into a better conditional mean-force
        estimate, \emph{provided} the birth--death does not corrupt
        $Y\mid X=x$. By \cref{prop:cond_invariance} a benign FR step preserves
        the conditional, so improved $x$-coverage translates into a better
        $\Fphat$ and hence a better $\Fhat$.
\end{enumerate}

\begin{figure}[t]
  \centering
  \includegraphics[width=0.95\textwidth]{fig05_mechanism.png}
  \caption{Mechanistic / early-time behaviour, estimated FR (best
    integrated-rule config) versus ABF only. (a) Early-time free-energy error
    (log scale): FR converges faster. (b) Region occupancy under FR relaxes
    from an unbalanced start toward the balanced $1/3$ split (dotted).
    (c) Cumulative barrier crossings are nearly identical --- the benefit here
    is redistribution, not a large increase in raw crossings. (d) FR score
    dispersion decays as $\phat_t\to q_t$.}
  \label{fig:mechanism}
\end{figure}

A notable feature of \cref{fig:mechanism}(c) is that FR does \emph{not} produce
many more barrier crossings than ABF in this example; the benefit comes from
\emph{reallocating} existing replicas toward starved values of $x$ rather than
from manufacturing rare transitions. Consistent with this, the event fractions
throughout (\cref{fig:convergence}, bottom right) are of order $10^{-5}$: in the
successful regimes FR is a \emph{weak} correction.

\subsection{FR does not visibly corrupt \texorpdfstring{$Y\mid X=x$}{Y given X}}
The central risk of a configuration-resampling step is that it distorts the
conditional law $Y\mid X=x$ that ABF relies on. \Cref{fig:conditional} compares
the conditional fidelity of all methods at the final time. The FR variants are
\emph{not} worse than ABF only; near the barrier ($x\in\{-0.5,0,0.5\}$), where
ABF-only under-samples, the FR variants are in fact slightly \emph{better} on
both the $\Ltwo$ and KL conditional metrics. This is direct evidence that, in
the successful regimes, the FR mechanism improves $x$-coverage without
corrupting the orthogonal sampling --- exactly the benign behaviour
\cref{prop:cond_invariance} requires.

\begin{figure}[t]
  \centering
  \includegraphics[width=0.95\textwidth]{fig06_conditional_y_given_x.png}
  \caption{Conditional $Y\mid X=x$ fidelity at $T$, by $x$-bin: (a) $\Ltwo$
    distance and (b) KL divergence of the empirical $\phat(y\mid x)$ to the
    reference $\pref(y\mid x)$. FR does not inflate the conditional error; near
    the barrier it slightly reduces it. \emph{Limitation:} the conditional CSV
    identifies runs only by (method, target, seed) and is aggregated over
    configurations, so this figure cannot attribute differences to specific
    good/bad hyperparameters (\cref{sec:discussion}).}
  \label{fig:conditional}
\end{figure}

\begin{remark}[Conditional diagnostics are not configuration-resolved]
\label{rem:cond_limit}
The current conditional-diagnostics output records only (method, target type,
seed) per row, not the FR hyperparameters. It therefore supports the
\emph{aggregate} statement that successful FR settings do not visibly corrupt
$Y\mid X=x$, but it cannot isolate a specific ``good'' versus ``bad''
configuration's effect on the conditional. A targeted good-versus-bad attribution
would require configuration-resolved conditional output; we list this as a
concrete improvement in \cref{sec:discussion}.
\end{remark}

\subsection{What determines a good FR? (Q3)}
\label{sec:q3}
The results let us name the factors that govern whether FR helps, and how.

\paragraph{Target quality is the practical bottleneck.}
The estimated target captures $\EstIntegImpPct\%$ of the integrated-error
reduction against the oracle's $\OracIntegImpPct\%$ (\cref{tab:main_integrated}).
Since the estimated and oracle targets share the \emph{same} mechanism and
differ only in whether the target free energy is the online estimate or
$\Fref$, this gap is attributable to the target estimate, not the birth--death.
Improving the online target free-energy estimate (its accuracy and its
stability under the EMA) is therefore the most direct route to a better
practical FR.

\paragraph{FR strength $\gamma$ trades speed against final accuracy.}
\Cref{fig:hyper} (estimated target) shows the integrated error decreasing as
$\gamma$ grows --- a stronger correction converges faster --- which is why the
integrated rule selects the aggressive $\gamma=\EstIntGamma$. But the same
aggressive setting does \emph{not} minimise the final error (\cref{sec:selection}):
the final-accuracy rule prefers the gentle $\gamma=\EstFinGamma$. There is thus
an explicit speed-versus-accuracy trade-off in $\gamma$.

\paragraph{Bandwidth $\eta$ and burn-in $b$.}
The marginal bandwidth $\eta$ sets the resolution at which over/under-population
is detected; intermediate values are robust in \cref{fig:hyper}. The burn-in $b$
controls when FR begins: starting FR early (small $b$) helps the
\emph{integrated} error here, because the largest imbalance is at early times,
whereas a later burn-in is safer for final accuracy. We caution against
universalising ``start early'': in this example ABF builds a usable target
quickly, so an early target is already informative; in harder problems a premature
target could misdirect the birth--death.

\paragraph{Event fraction and diversity.}
Across the successful configurations the event fractions are tiny
($\sim10^{-5}$) and the cap $f_{\max}$ is never the binding constraint, so the
production configurations are \emph{gentle} and particle diversity is preserved
(no collapse of ancestors). A more aggressive birth--death --- larger $\gamma$,
larger event fractions --- could in principle still harm accuracy by eroding
diversity or by acting on a noisy target; the present configurations simply do
not reach that regime.

\paragraph{Conditional fidelity.}
Finally, a good FR must preserve $Y\mid X=x$ (\cref{fig:conditional}). The
successful configurations do; this is a necessary condition for the coverage
gain to translate into a mean-force gain rather than a bias.

\begin{figure}[t]
  \centering
  \includegraphics[width=0.95\textwidth]{fig07_hyperparameter_heatmaps.png}
  \caption{Estimated-target FR hyperparameter sensitivity: median integrated
    $\Ltwo(F)$ over the grid (lower = faster convergence). Left: $\gamma$ versus
    $\eta$; right: $\gamma$ versus burn-in fraction. The integrated error
    improves with larger $\gamma$ and earlier burn-in in this example;
    intermediate $\eta$ is robust. These trends inform Q3 but should not be
    extrapolated beyond this regime.}
  \label{fig:hyper}
\end{figure}
